\chapter{Testy aplikacji}
Aplikacja powinna być testowalna w~każdej jej warstwie. W~zależności od warstwy testy podzielono na testy jednostkowe,
integracyjne i~testy akceptacyjne (\textit{End-to-End}). Testy jednostkowe sprawdzają poprawność logiki zawartej w~kodzie aplikacji,
testy integracyjne weryfikują poprawność współpracy między komponentami systemu, a~testy akceptacyjne sprawdzają
poprawność działania aplikacji z~perspektywy użytkownika końcowego. W~tym rozdziale przedstawiono podejście do testowania aplikacji wraz z~opisem narzędzi wykorzystanych do przeprowadzenia
testów. Zaprezentowano również przykłady testów jednostkowych, integracyjnych i~akceptacyjnych razem ze scenariuszami.

\section{Narzędzia}
Do implementacji testów wykorzystano frameworki i~biblioteki testowe, które pozwalają na automatyzację procesu testowania. W~zależności od rodzaju testów zastosowano różne narzędzia:
\begin{itemize}
\item \textbf{TUnit} --- framework do testów w~języku C\#.
\item \textbf{Microsoft Testing Platform} --- środowisko uruchomieniowe (runner) testów jednostkowych i~integracyjnych w~języku C\#.
\item \textbf{Microsoft Playwright} --- narzędzie do testów akceptacyjnych emulujące przeglądarkę internetową.
\item \textbf{Docker i~Testcontainers} --- narzędzia do konteneryzacji umożliwiające uruchamianie izolowanych kontenerów na potrzeby testów.
\end{itemize}
Ich dobór pozwala na automatyzację procesu testowania i~przyspieszenie weryfikacji poprawności działania systemu.

\section{Testy jednostkowe}
Testy jednostkowe obejmują warstwę domeny aplikacji, gdzie znajdują się agregaty zawierające logikę biznesową.
Testowano poprawność wytwarzania danych, interakcji między obiektami, walidacji i~zachowania samych obiektów.
Testy jednostkowe zostały zaimplementowane w~języku C\# z~wykorzystaniem frameworka TUnit. W~sekcji tej przedstawiono
kilka scenariuszy testowych, które zostały stworzone na potrzeby testowania systemu. 

Każdy test tworzy encje bezpośrednio w~pamięci i~wywołuje metody agregatu, bez bazy danych i~bez warstwy HTTP.
Testy są asynchroniczne. Testy oznaczone atrybutem \texttt{[Arguments]} są uruchamiane osobno dla każdego zestawu
podanych argumentów.

\subsection{Scenariusz 1: domyślne role nowego projektu}
Utworzenie projektu tworzy trzy role: \textit{Administrator}, \textit{Członek} i~\textit{Obserwator}. Uprawnienia ról są
wyliczane z~katalogu uprawnień, czyli ze wszystkich par zasób i~akcja. Scenariusz sprawdza, czy rola \textit{Członek}
dostaje dokładnie te uprawnienia, które przewiduje model uprawnień. Opis scenariusza zawiera tabela~\ref{tab:ut-role}.

\begin{table}[htb]
    \centering
    \caption{Scenariusz testowy: domyślne role nowego projektu}
    \label{tab:ut-role}
    \small
    \begin{tabularx}{\linewidth}{@{}>{\raggedright\arraybackslash}p{3.5cm} >{\raggedright\arraybackslash}X@{}}
        \hline
        Test & \texttt{ProjectAggregateTests.\allowbreak{}Create\_\allowbreak{}Member\_\allowbreak{}HasViewEverywhere\_\allowbreak{}FullManagement\_\allowbreak{}AndRepositoryCreate} \\ \hline
        Warunki początkowe & Katalog uprawnień zawiera wszystkie pary zasób i~akcja (6 zasobów, 4 akcje) \\ \hline
        Kroki & Wywołanie \texttt{ProjectAggregate.Create} z~nazwą, opisem, identyfikatorem właściciela i~katalogiem uprawnień \\ \hline
        Oczekiwany wynik & Rola \textit{Członek} ma uprawnienie \textit{View} do każdego zasobu, wszystkie akcje na zasobie \textit{Management} oraz uprawnienie \textit{Create} na zasobie \textit{Repositories}. Nie ma żadnych innych uprawnień \\ \hline
    \end{tabularx}
\end{table}

Implementację testu przedstawiono na listingu~\ref{lst:ut-role}. Oczekiwany zbiór uprawnień jest budowany z~typów
wyliczeniowych \texttt{ResourceType} i~\texttt{ActionType}, dzięki czemu po dodaniu nowego zasobu test nadal sprawdza
poprawną regułę.

\begin{lstlisting}[basicstyle=\footnotesize\ttfamily, captionpos=t, caption={Test uprawnień domyślnej roli \textit{Członek}}, label={lst:ut-role}]
[Test]
public async Task Create_Member_HasViewEverywhere_FullManagement_AndRepositoryCreate()
{
    var project = ProjectAggregate.Create("Nazwa", "Opis", Guid.NewGuid(), _catalog);

    var permissions = project.Role(Member).FeatureActions.Select(x => x.Key());
    var expected = Enum.GetValues<ResourceType>().Select(r => $"{r}.{ActionType.View}")
        .Concat(Enum.GetValues<ActionType>().Select(a => $"{ResourceType.Management}.{a}"))
        .Append($"{ResourceType.Repositories}.{ActionType.Create}")
        .Distinct();

    await Assert.That(permissions).IsEquivalentTo(expected);
}
\end{lstlisting}

\subsection{Scenariusz 2: unikalna nazwa roli w~projekcie}
Nazwa roli musi być unikalna w~obrębie projektu. Porównanie ignoruje wielkość liter oraz białe znaki na początku i~końcu
nazwy, ponieważ użytkownik nie odróżni ról \textit{Członek} i~\textit{członek}. Opis scenariusza zawiera tabela~\ref{tab:ut-nazwa-roli}.

\begin{table}[htb]
    \centering
    \caption{Scenariusz testowy: unikalna nazwa roli}
    \label{tab:ut-nazwa-roli}
    \small
    \begin{tabularx}{\linewidth}{@{}>{\raggedright\arraybackslash}p{3.5cm} >{\raggedright\arraybackslash}X@{}}
        \hline
        Test & \texttt{ProjectAggregateTests.\allowbreak{}AddRole\_\allowbreak{}DuplicateNameIgnoringCaseAndWhitespace\_\allowbreak{}ThrowsConflict} \\ \hline
        Warunki początkowe & Projekt z~trzema domyślnymi rolami \\ \hline
        Dane wejściowe & \texttt{"Członek"}, \texttt{"członek"}, \texttt{"  ADMINISTRATOR  "} \\ \hline
        Kroki & Wywołanie \texttt{AddRole} z~nazwą istniejącej roli \\ \hline
        Oczekiwany wynik & Wyjątek \texttt{ProjectAccessException} z~kodem \texttt{409 Conflict}. Projekt nadal ma trzy role \\ \hline
    \end{tabularx}
\end{table}

Test z~listingu~\ref{lst:ut-nazwa-roli} sprawdza także, czy nieudana operacja nie zmieniła stanu agregatu.

\begin{lstlisting}[basicstyle=\footnotesize\ttfamily, captionpos=t, caption={Test odrzucenia zduplikowanej nazwy roli}, label={lst:ut-nazwa-roli}]
[Test]
[Arguments("Członek")]
[Arguments("członek")]
[Arguments("  ADMINISTRATOR  ")]
public async Task AddRole_DuplicateNameIgnoringCaseAndWhitespace_ThrowsConflict(string name)
{
    var project = TestEntities.Project(catalog: _catalog);
    var aggregate = new ProjectAggregate(project);

    var exception = await Assert.That(() => aggregate.AddRole(name, false, []))
        .Throws<ProjectAccessException>();

    await Assert.That(exception!.StatusCode).IsEqualTo(HttpStatusCode.Conflict);
    await Assert.That(project.ProjectRoles.Count).IsEqualTo(3);
}
\end{lstlisting}

\subsection{Scenariusz 3: kolejne numery zadań w~projekcie}
Każde zadanie ma przyjazny numer (\textit{FriendlyId}) wyświetlany użytkownikowi, np. \texttt{\#3}. Numery są nadawane
kolejno w~obrębie projektu na podstawie licznika \texttt{LastTaskFriendlyId}. Scenariusz sprawdza nadawanie numerów oraz to,
że odrzucone zadanie nie zużywa numeru. Opis scenariusza zawiera tabela~\ref{tab:ut-numery}.

\begin{table}[htb]
    \centering
    \caption{Scenariusz testowy: kolejne numery zadań}
    \label{tab:ut-numery}
    \small
    \begin{tabularx}{\linewidth}{@{}>{\raggedright\arraybackslash}p{3.5cm} >{\raggedright\arraybackslash}X@{}}
        \hline
        Testy & \texttt{TaskAggregateTests.\allowbreak{}Create\_\allowbreak{}AssignsConsecutiveFriendlyIds}, \texttt{TaskAggregateTests.\allowbreak{}Create\_\allowbreak{}SprintFromAnotherProject\_\allowbreak{}Throws\_\allowbreak{}AndDoesNotConsumeFriendlyId} \\ \hline
        Warunki początkowe & Nowy projekt z~licznikiem zadań równym 0 \\ \hline
        Kroki & 1. Utworzenie trzech zadań w~projekcie. \newline 2. W~osobnym teście: próba utworzenia zadania w~sprincie należącym do innego projektu \\ \hline
        Oczekiwany wynik & 1. Zadania mają numery 1, 2 i~3, a~licznik projektu wynosi 3. \newline 2. Wyjątek \texttt{InvalidOperationException}, a~licznik projektu pozostaje równy 0 \\ \hline
    \end{tabularx}
\end{table}

Oba testy przedstawiono na listingu~\ref{lst:ut-numery}.

\begin{lstlisting}[basicstyle=\footnotesize\ttfamily, captionpos=t, caption={Testy nadawania numerów zadań}, label={lst:ut-numery}]
[Test]
public async Task Create_AssignsConsecutiveFriendlyIds()
{
    var first = TaskAggregate.Create(_project, "A", "", TaskPriority.Normal, TaskStatus.Todo, null, null);
    var second = TaskAggregate.Create(_project, "B", "", TaskPriority.Normal, TaskStatus.Todo, null, null);
    var third = TaskAggregate.Create(_project, "C", "", TaskPriority.Normal, TaskStatus.Todo, null, null);

    await Assert.That(new[] { first.FriendlyId, second.FriendlyId, third.FriendlyId })
        .IsEquivalentTo(new[] { 1L, 2L, 3L });
    await Assert.That(_project.LastTaskFriendlyId).IsEqualTo(3L);
}

[Test]
public async Task Create_SprintFromAnotherProject_Throws_AndDoesNotConsumeFriendlyId()
{
    var foreignSprint = TestEntities.Sprint(Guid.NewGuid());

    await Assert.That(() =>
            TaskAggregate.Create(_project, "A", "", TaskPriority.Normal, TaskStatus.Todo, foreignSprint, null))
        .Throws<InvalidOperationException>();

    await Assert.That(_project.LastTaskFriendlyId).IsEqualTo(0L);
}
\end{lstlisting}

\subsection{Scenariusz 4: data ukończenia zadania}
Zmiana statusu zadania na \textit{Done} zapisuje datę ukończenia. Przywrócenie zadania do statusu \textit{Todo} lub
\textit{InProgress} usuwa tę datę, żeby statystyki sprintu nie liczyły go jako ukończonego. Opis scenariusza zawiera
tabela~\ref{tab:ut-status}.

\begin{table}[htb]
    \centering
    \caption{Scenariusz testowy: data ukończenia zadania}
    \label{tab:ut-status}
    \small
    \begin{tabularx}{\linewidth}{@{}>{\raggedright\arraybackslash}p{3.5cm} >{\raggedright\arraybackslash}X@{}}
        \hline
        Testy & \texttt{TaskAggregateTests.\allowbreak{}SetStatus\_\allowbreak{}ToDone\_\allowbreak{}SetsFinishedDate}, \texttt{TaskAggregateTests.\allowbreak{}SetStatus\_\allowbreak{}FromDoneToOther\_\allowbreak{}ClearsFinishedDate} \\ \hline
        Warunki początkowe & 1. Zadanie w~statusie \textit{InProgress}. \newline 2. Zadanie w~statusie \textit{Done} z~datą ukończenia 1 stycznia 2026 \\ \hline
        Dane wejściowe & Dla drugiego testu: \textit{Todo} i~\textit{InProgress} \\ \hline
        Kroki & Wywołanie \texttt{SetStatus} z~nowym statusem \\ \hline
        Oczekiwany wynik & 1. Status \textit{Done}, data ukończenia jest ustawiona. \newline 2. Nowy status, brak daty ukończenia \\ \hline
    \end{tabularx}
\end{table}

Implementację testów przedstawiono na listingu~\ref{lst:ut-status}.

\begin{lstlisting}[basicstyle=\footnotesize\ttfamily, captionpos=t, caption={Testy daty ukończenia przy zmianie statusu zadania}, label={lst:ut-status}]
[Test]
public async Task SetStatus_ToDone_SetsFinishedDate()
{
    var task = TestEntities.Task(_project.Id, TaskStatus.InProgress);

    new TaskAggregate(task).SetStatus(TaskStatus.Done);

    await Assert.That(task.Status).IsEqualTo(TaskStatus.Done);
    await Assert.That(task.FinishedDate).IsNotNull();
}

[Test]
[Arguments(TaskStatus.Todo)]
[Arguments(TaskStatus.InProgress)]
public async Task SetStatus_FromDoneToOther_ClearsFinishedDate(TaskStatus status)
{
    var task = TestEntities.Task(_project.Id, TaskStatus.Done);
    task.FinishedDate = new DateOnly(2026, 1, 1);

    new TaskAggregate(task).SetStatus(status);

    await Assert.That(task.Status).IsEqualTo(status);
    await Assert.That(task.FinishedDate).IsNull();
}
\end{lstlisting}

\subsection{Scenariusz 5: jednokrotna odpowiedź na zaproszenie}
Na zaproszenie do projektu można odpowiedzieć tylko raz. Przyjęcie zaproszenia podnosi zdarzenie domenowe
\texttt{InvitationAcceptedEvent}, którego obsługa dodaje użytkownika do projektu. Druga odpowiedź mogłaby dodać członka
po wcześniejszej odmowie, dlatego agregat ją odrzuca. Test sprawdza wszystkie cztery kombinacje pierwszej i~drugiej
odpowiedzi. Opis scenariusza zawiera tabela~\ref{tab:ut-zaproszenie}.

\begin{table}[htb]
    \centering
    \caption{Scenariusz testowy: jednokrotna odpowiedź na zaproszenie}
    \label{tab:ut-zaproszenie}
    \small
    \begin{tabularx}{\linewidth}{@{}>{\raggedright\arraybackslash}p{3.5cm} >{\raggedright\arraybackslash}X@{}}
        \hline
        Test & \texttt{ProjectInvitationAggregateTests.\allowbreak{}AnsweringTwice\_\allowbreak{}ThrowsConflict\_\allowbreak{}AndKeepsFirstAnswer} \\ \hline
        Warunki początkowe & Aktywne zaproszenie do projektu \\ \hline
        Dane wejściowe & Pary (pierwsza odpowiedź, druga odpowiedź): (przyjęcie, przyjęcie), (przyjęcie, odmowa), (odmowa, przyjęcie), (odmowa, odmowa) \\ \hline
        Kroki & 1. Pierwsza odpowiedź na zaproszenie. \newline 2. Druga odpowiedź na to samo zaproszenie \\ \hline
        Oczekiwany wynik & Druga odpowiedź kończy się wyjątkiem \texttt{ProjectAccessException} z~kodem \texttt{409 Conflict}. Zaproszenie zachowuje pierwszą odpowiedź, a~liczba zdarzeń domenowych się nie zmienia \\ \hline
    \end{tabularx}
\end{table}

Implementację testu przedstawiono na listingu~\ref{lst:ut-zaproszenie}.

\begin{lstlisting}[basicstyle=\footnotesize\ttfamily, captionpos=t, caption={Test ponownej odpowiedzi na zaproszenie}, label={lst:ut-zaproszenie}]
[Test]
[Arguments(true, true)]
[Arguments(true, false)]
[Arguments(false, true)]
[Arguments(false, false)]
public async Task AnsweringTwice_ThrowsConflict_AndKeepsFirstAnswer(bool firstAccept, bool secondAccept)
{
    var invitation = NewInvitation();
    var aggregate = new ProjectInvitationAggregate(invitation);
    Answer(aggregate, firstAccept);
    var eventsAfterFirstAnswer = aggregate.DomainEvents.Count;

    var exception = await Assert.That(() => Answer(aggregate, secondAccept))
        .Throws<ProjectAccessException>();

    await Assert.That(exception!.StatusCode).IsEqualTo(HttpStatusCode.Conflict);
    await Assert.That(invitation.UserAccepted).IsEqualTo(firstAccept);
    await Assert.That(aggregate.DomainEvents.Count).IsEqualTo(eventsAfterFirstAnswer);
}
\end{lstlisting}

\section{Testy integracyjne}
Testy integracyjne pokrywają kilka warstw aplikacji, w~tym warstwę domeny, infrastruktury i~warstwę aplikacji.
Mają one na celu weryfikację, czy system poprawnie komunikuje się z~bazą danych i~zewnętrznymi usługami oraz czy prawidłowo przetwarza dane w~całym cyklu życia żądania. Testy integracyjne są uruchamiane w~środowisku zbliżonym do produkcyjnego, często z~użyciem kontenerów \textit{Docker}, aby zapewnić spójność środowiska testowego.
W~tym przypadku wykorzystano bibliotekę \textit{Testcontainers}, która tworzy tymczasowe kontenery \textit{Docker} na czas przeprowadzania testów.
Weryfikowano tutaj, czy aplikacja poprawnie zapisuje i~odczytuje dane z~bazy danych, czy wysyła powiadomienia e-mail do właściwych odbiorców, czy poprawnie
obsługuje repozytoria \textit{Git} oraz czy prawidłowo rotuje token odświeżania (\textit{RefreshToken}).

Testy zaimplementowano w~języku C\# z~wykorzystaniem frameworka TUnit oraz klasy \texttt{WebApplicationFactory}
z~biblioteki \textit{Microsoft.AspNetCore.Mvc.Testing}. Klasa ta uruchamia całe API w~procesie testów, razem z~migracjami
bazy danych i~danymi początkowymi. Bazą danych jest PostgreSQL w~kontenerze tworzonym przez \textit{Testcontainers}.
API nasłuchuje na prawdziwym porcie, dzięki czemu do testów repozytoriów można użyć klienta \textit{Git} z~wiersza poleceń.
Usługa wysyłki wiadomości e-mail jest zastąpiona implementacją testową, która zapisuje wiadomości w~pamięci.
Jedna instancja API i~jeden kontener są współdzielone przez wszystkie testy. Każdy test rejestruje własnych użytkowników
i~tworzy własne projekty, więc testy mogą działać równolegle i~nie trzeba czyścić bazy danych między nimi.
Konfigurację środowiska testowego przedstawiono na listingu~\ref{lst:it-fabryka}.

\begin{lstlisting}[basicstyle=\footnotesize\ttfamily, captionpos=t, caption={Fragment klasy uruchamiającej API na potrzeby testów integracyjnych}, label={lst:it-fabryka}]
public sealed class ApiFactory : WebApplicationFactory<Program>, IAsyncInitializer
{
    private readonly PostgreSqlContainer _postgres = new PostgreSqlBuilder()
        .WithImage("postgres:18-alpine")
        .WithDatabase("ProjectCowork")
        .Build();

    public FakeMailService Mail { get; } = new();

    public ApiFactory()
    {
        UseKestrel(0);
    }

    protected override void ConfigureWebHost(IWebHostBuilder builder)
    {
        builder.UseEnvironment("Testing");
        builder.UseSetting("ConnectionStrings:DefaultConnection", _postgres.GetConnectionString());

        builder.ConfigureTestServices(services =>
        {
            services.RemoveAll<IMailService>();
            services.AddSingleton<IMailService>(Mail);
        });
    }
}
\end{lstlisting}

\subsection{Scenariusz 1: rotacja tokena odświeżania}
Token odświeżania jest przechowywany w~ciasteczku \texttt{HttpOnly}. Przy każdym odświeżeniu
sesji API wydaje nowy token i~unieważnia poprzedni. Dzięki temu skradziony token można wykorzystać najwyżej raz.
Scenariusz sprawdza, czy po odświeżeniu sesji stary token przestaje działać. Opis scenariusza zawiera tabela~\ref{tab:it-rotacja}.

\begin{table}[htb]
    \centering
    \caption{Scenariusz testowy: rotacja tokena odświeżania}
    \label{tab:it-rotacja}
    \small
    \begin{tabularx}{\linewidth}{@{}>{\raggedright\arraybackslash}p{3.5cm} >{\raggedright\arraybackslash}X@{}}
        \hline
        Test & \texttt{AuthTests.\allowbreak{}Refresh\_\allowbreak{}RotatesToken\_\allowbreak{}OldTokenCannotBeReused} \\ \hline
        Warunki początkowe & Nowo zarejestrowany użytkownik z~tokenem odświeżania z~odpowiedzi na rejestrację \\ \hline
        Kroki & 1. Wywołanie \texttt{POST /api/auth/token/refresh} z~tokenem z~rejestracji. \newline 2. Ponowne wywołanie z~tym samym tokenem \\ \hline
        Oczekiwany wynik & 1. Kod \texttt{200 OK} i~nowy token. \newline 2. Kod \texttt{401 Unauthorized} \\ \hline
    \end{tabularx}
\end{table}

Implementację testu przedstawiono na listingu~\ref{lst:it-rotacja}.

\begin{lstlisting}[basicstyle=\footnotesize\ttfamily, captionpos=t, caption={Test rotacji tokena odświeżania}, label={lst:it-rotacja}]
[Test]
public async Task Refresh_RotatesToken_OldTokenCannotBeReused()
{
    var user = await RegisterAsync();
    var first = await PostWithRefreshCookieAsync("/api/auth/token/refresh", user.RefreshToken);
    await EnsureStatusAsync(first, HttpStatusCode.OK);

    var replay = await PostWithRefreshCookieAsync("/api/auth/token/refresh", user.RefreshToken);

    await Assert.That(replay.StatusCode).IsEqualTo(HttpStatusCode.Unauthorized);
}
\end{lstlisting}

\subsection{Scenariusz 2: uprawnienia do tworzenia zadań}
Uprawnienia są sprawdzane w~obrębie projektu na podstawie roli członka. Scenariusz sprawdza, czy reguły z~domyślnych ról
działają po przejściu przez całą ścieżkę żądania: uwierzytelnienie tokenem JWT, sprawdzenie uprawnień w~bazie danych
i~obsługę wyjątków w~API. Wszystkie przypadki testowe korzystają z~jednego projektu, do którego dołączono użytkowników
z~każdą z~domyślnych ról. Opis scenariusza zawiera tabela~\ref{tab:it-uprawnienia}.

\begin{table}[htb]
    \centering
    \caption{Scenariusz testowy: uprawnienia do tworzenia zadań}
    \label{tab:it-uprawnienia}
    \small
    \begin{tabularx}{\linewidth}{@{}>{\raggedright\arraybackslash}p{3.5cm} >{\raggedright\arraybackslash}X@{}}
        \hline
        Test & \texttt{PermissionMatrixTests.\allowbreak{}AddTask\_\allowbreak{}RequiresManagementCreate} \\ \hline
        Warunki początkowe & Projekt z~właścicielem oraz członkami w~rolach \textit{Administrator}, \textit{Członek} i~\textit{Obserwator} \\ \hline
        Dane wejściowe & Właściciel, Administrator, Członek, Obserwator, użytkownik spoza projektu, użytkownik niezalogowany \\ \hline
        Kroki & Wywołanie \texttt{POST /api/tasks} w~imieniu każdego z~użytkowników \\ \hline
        Oczekiwany wynik & Właściciel, Administrator i~Członek: \texttt{201 Created}. Obserwator i~użytkownik spoza projektu: \texttt{403 Forbidden}. Użytkownik niezalogowany: \texttt{401 Unauthorized} \\ \hline
    \end{tabularx}
\end{table}

Implementację testu przedstawiono na listingu~\ref{lst:it-uprawnienia}. W~ten sam sposób przetestowano pozostałe
endpointy, m.in. odczyt projektu, edycję projektu, tworzenie sprintów i~zarządzanie rolami.

\begin{lstlisting}[basicstyle=\footnotesize\ttfamily, captionpos=t, caption={Test uprawnień do tworzenia zadań}, label={lst:it-uprawnienia}]
[Test]
[Arguments(Actor.Owner, HttpStatusCode.Created)]
[Arguments(Actor.Administrator, HttpStatusCode.Created)]
[Arguments(Actor.Member, HttpStatusCode.Created)]
[Arguments(Actor.Observer, HttpStatusCode.Forbidden)]
[Arguments(Actor.Outsider, HttpStatusCode.Forbidden)]
[Arguments(Actor.Anonymous, HttpStatusCode.Unauthorized)]
public async Task AddTask_RequiresManagementCreate(Actor actor, HttpStatusCode expected)
{
    var scenario = await GetScenarioAsync();

    var response = await PostTaskAsync(ClientFor(scenario, actor), scenario.Project.Id, Unique("Zadanie"));

    await Assert.That(response.StatusCode).IsEqualTo(expected);
}
\end{lstlisting}

\subsection{Scenariusz 3: dołączenie do projektu przez zaproszenie}
Użytkownik dołącza do projektu, akceptując zaproszenie. Akceptacja podnosi zdarzenie domenowe, którego obsługa dodaje
użytkownika do projektu z~rolą wskazaną w~zaproszeniu. Scenariusz sprawdza, czy zdarzenie zostało obsłużone, a~zmiany
zapisane w~bazie danych. Opis scenariusza zawiera tabela~\ref{tab:it-zaproszenie}.

\begin{table}[htb]
    \centering
    \caption{Scenariusz testowy: dołączenie do projektu przez zaproszenie}
    \label{tab:it-zaproszenie}
    \small
    \begin{tabularx}{\linewidth}{@{}>{\raggedright\arraybackslash}p{3.5cm} >{\raggedright\arraybackslash}X@{}}
        \hline
        Test & \texttt{InvitationTests.\allowbreak{}Accept\_\allowbreak{}WithRole\_\allowbreak{}AddsMemberWithThatRole} \\ \hline
        Warunki początkowe & Dwóch zarejestrowanych użytkowników. Pierwszy z~nich jest właścicielem projektu \\ \hline
        Kroki & 1. Właściciel zaprasza drugiego użytkownika z~rolą \textit{Obserwator}. \newline 2. Zaproszony pobiera listę swoich zaproszeń i~akceptuje zaproszenie. \newline 3. Właściciel pobiera listę członków projektu \\ \hline
        Oczekiwany wynik & Lista zawiera dwie osoby. Pierwszą jest właściciel, oznaczony jako właściciel projektu. Drugą jest zaproszony użytkownik z~rolą \textit{Obserwator} \\ \hline
    \end{tabularx}
\end{table}

Implementację testu przedstawiono na listingu~\ref{lst:it-zaproszenie}. Zaproszenie i~jego akceptację wykonuje metoda
pomocnicza \texttt{AddMemberAsync}, która wywołuje kolejno endpointy zaproszenia, listy zaproszeń i~akceptacji.

\begin{lstlisting}[basicstyle=\footnotesize\ttfamily, captionpos=t, caption={Test dołączenia do projektu z~wybraną rolą}, label={lst:it-zaproszenie}]
[Test]
public async Task Accept_WithRole_AddsMemberWithThatRole()
{
    var owner = await RegisterAsync();
    var invitee = await RegisterAsync();
    var project = await CreateProjectAsync(owner);

    await AddMemberAsync(owner, project.Id, invitee, ObserverRole);

    var members = await GetListAsync<ProjectMemberDto>(owner.Client, $"/api/projects/{project.Id}/members");
    await Assert.That(members.Count).IsEqualTo(2);
    await Assert.That(members[0].IsOwner).IsTrue();
    await Assert.That(members[0].UserId).IsEqualTo(owner.Id);

    var member = members.Single(x => x.UserId == invitee.Id);
    await Assert.That(member.RoleName).IsEqualTo(ObserverRole);
    await Assert.That(member.IsOwner).IsFalse();
}
\end{lstlisting}

\subsection{Scenariusz 4: usunięcie sprintu}
Usunięcie sprintu nie może usuwać jego zadań. Zadania wracają do backlogu, czyli tracą przypisanie do sprintu. Odpowiada
za to obsługa zdarzenia domenowego \texttt{SprintDeletedEvent}, a~zmiany są zapisywane razem z~usunięciem sprintu.
Scenariusz sprawdza także, czy zadania z~innego sprintu pozostają bez zmian. Opis scenariusza zawiera tabela~\ref{tab:it-sprint}.

\begin{table}[htb]
    \centering
    \caption{Scenariusz testowy: usunięcie sprintu}
    \label{tab:it-sprint}
    \small
    \begin{tabularx}{\linewidth}{@{}>{\raggedright\arraybackslash}p{3.5cm} >{\raggedright\arraybackslash}X@{}}
        \hline
        Test & \texttt{SprintAndTaskTests.\allowbreak{}DeleteSprint\_\allowbreak{}MovesItsTasksToBacklog} \\ \hline
        Warunki początkowe & Projekt z~dwoma sprintami. W~pierwszym sprincie są dwa zadania, w~drugim jedno \\ \hline
        Kroki & 1. Wywołanie \texttt{DELETE /api/sprints/\{id\}} dla pierwszego sprintu. \newline 2. Pobranie zadań i~sprintów projektu \\ \hline
        Oczekiwany wynik & Kod \texttt{204 No Content}. Projekt nadal ma trzy zadania. Dwa zadania z~usuniętego sprintu nie są przypisane do żadnego sprintu, a~zadanie z~drugiego sprintu jest w~nim nadal. Projekt ma jeden sprint \\ \hline
    \end{tabularx}
\end{table}

Implementację testu przedstawiono na listingu~\ref{lst:it-sprint}.

\begin{lstlisting}[basicstyle=\footnotesize\ttfamily, captionpos=t, caption={Test przeniesienia zadań do backlogu po usunięciu sprintu}, label={lst:it-sprint}]
[Test]
public async Task DeleteSprint_MovesItsTasksToBacklog()
{
    var owner = await RegisterAsync();
    var project = await CreateProjectAsync(owner);
    var sprint = await CreateSprintAsync(owner, project.Id);
    var otherSprint = await CreateSprintAsync(owner, project.Id);
    var first = await CreateTaskAsync(owner, project.Id, sprintId: sprint.Id);
    var second = await CreateTaskAsync(owner, project.Id, sprintId: sprint.Id);
    var untouched = await CreateTaskAsync(owner, project.Id, sprintId: otherSprint.Id);

    var response = await owner.Client.DeleteAsync($"/api/sprints/{sprint.Id}");

    await Assert.That(response.StatusCode).IsEqualTo(HttpStatusCode.NoContent);

    var tasks = (await GetTasksAsync(owner, project.Id)).ToDictionary(x => x.Id);
    await Assert.That(tasks.Count).IsEqualTo(3);
    await Assert.That(tasks[first.Id].SprintId).IsNull();
    await Assert.That(tasks[second.Id].SprintId).IsNull();
    await Assert.That(tasks[untouched.Id].SprintId).IsEqualTo(otherSprint.Id);

    var sprints = await GetListAsync<SprintDto>(owner.Client, $"/api/sprints/project/{project.Id}");
    await Assert.That(sprints.Select(x => x.Id)).IsEquivalentTo(new[] { otherSprint.Id });
}
\end{lstlisting}

\subsection{Scenariusz 5: historia commitów po wypchnięciu zmian}
Scenariusz sprawdza współpracę API z~serwerem \textit{Git}. Test tworzy repozytorium przez API, a~następnie klonuje je,
tworzy dwa commity i~wypycha je poleceniem \texttt{git push}, tak jak zrobiłby to użytkownik. Historia commitów
odczytana przez API musi odpowiadać temu, co zostało wypchnięte. Opis scenariusza zawiera tabela~\ref{tab:it-commity}.

\begin{table}[htb]
    \centering
    \caption{Scenariusz testowy: historia commitów po wypchnięciu zmian}
    \label{tab:it-commity}
    \small
    \begin{tabularx}{\linewidth}{@{}>{\raggedright\arraybackslash}p{3.5cm} >{\raggedright\arraybackslash}X@{}}
        \hline
        Test & \texttt{RepositoryTests.\allowbreak{}Commits\_\allowbreak{}NewestFirst\_\allowbreak{}WithPaging} \\ \hline
        Warunki początkowe & Repozytorium utworzone przez API, do którego wypchnięto commity \textit{Pierwszy commit} i~\textit{Drugi commit} \\ \hline
        Kroki & 1. Pobranie wszystkich commitów gałęzi \texttt{main}. \newline 2. Pobranie pierwszej i~drugiej strony z~jednym commitem na stronę \\ \hline
        Oczekiwany wynik & 1. Dwa commity, od najnowszego, z~autorem i~40-znakowym skrótem SHA. Brak kolejnej strony. \newline 2. Pierwsza strona zawiera \textit{Drugi commit} i~informację o~kolejnej stronie. Druga strona zawiera \textit{Pierwszy commit} i~jest ostatnia \\ \hline
    \end{tabularx}
\end{table}

Przygotowanie repozytorium przedstawiono na listingu~\ref{lst:it-push}, a~sam test na listingu~\ref{lst:it-commity}.
Klasa \texttt{GitCli} uruchamia klienta \textit{Git} z~odciętą konfiguracją użytkownika, żeby wynik testu nie zależał
od ustawień komputera, na którym jest uruchamiany.

\begin{lstlisting}[basicstyle=\footnotesize\ttfamily, captionpos=t, caption={Przygotowanie repozytorium z~dwoma wypchniętymi commitami}, label={lst:it-push}]
var owner = await RegisterAsync();
var project = await CreateProjectAsync(owner);
var repository = await CreateRepositoryAsync(owner, project.Id);

var work = await git.CloneAsync(repository.RepoUrl, "work");
await git.RunAsync(work, "checkout", "-b", "main");

await GitCli.WriteFileAsync(work, "README.md", "# Projekt\n");
await git.CommitAllAsync(work, "Pierwszy commit");

await GitCli.WriteFileAsync(work, "src/app.txt", "Witaj świecie\n");
await File.WriteAllBytesAsync(Path.Combine(work, "logo.bin"), [0, 1, 2, 0, 255]);
await git.CommitAllAsync(work, "Drugi commit");

await git.RunAsync(work, "push", "origin", "main");
\end{lstlisting}

\begin{lstlisting}[basicstyle=\footnotesize\ttfamily, captionpos=t, caption={Test historii commitów i~stronicowania}, label={lst:it-commity}]
[Test]
public async Task Commits_NewestFirst_WithPaging()
{
    using var git = new GitCli();
    var (owner, _, repository) = await CreatePushedRepositoryAsync(git);

    var all = await owner.Client.GetFromJsonAsync<CommitsPageDto>(
        $"/api/repository/{repository.Id}/commits?branch=main");
    var firstPage = await owner.Client.GetFromJsonAsync<CommitsPageDto>(
        $"/api/repository/{repository.Id}/commits?branch=main&page=1&pageSize=1");
    var secondPage = await owner.Client.GetFromJsonAsync<CommitsPageDto>(
        $"/api/repository/{repository.Id}/commits?branch=main&page=2&pageSize=1");

    await Assert.That(all!.Items.Select(x => x.Message))
        .IsEquivalentTo(["Drugi commit", "Pierwszy commit"], CollectionOrdering.Matching);
    await Assert.That(all.HasMore).IsFalse();
    await Assert.That(all.Items[0].AuthorName).IsEqualTo("Tester");
    await Assert.That(all.Items[0].Sha.Length).IsEqualTo(40);

    await Assert.That(firstPage!.Items.Single().Message).IsEqualTo("Drugi commit");
    await Assert.That(firstPage.HasMore).IsTrue();
    await Assert.That(secondPage!.Items.Single().Message).IsEqualTo("Pierwszy commit");
    await Assert.That(secondPage.HasMore).IsFalse();
}
\end{lstlisting}

\section{Testy akceptacyjne}
Testy akceptacyjne, w~przeciwieństwie do testów jednostkowych i~integracyjnych, koncentrują się na weryfikacji, czy system spełnia wymagania biznesowe i~oczekiwania użytkowników końcowych. 
Testy te są zwykle przeprowadzane z~perspektywy użytkownika i~obejmują pełne scenariusze użycia aplikacji, od interakcji z~interfejsem użytkownika po weryfikację wyników w~systemie.
Zostały one przeprowadzone w~aplikacji klienckiej, za pomocą której użytkownicy wchodzą w~interakcję z~systemem.
Testy wykonano z~użyciem narzędzia \textit{Playwright}, które umożliwia automatyzację przeglądarki internetowej i~symulowanie działań użytkownika w~aplikacji webowej.

Testy akceptacyjne zaimplementowano w~języku \textit{TypeScript} z~wykorzystaniem biblioteki \textit{Playwright Test}
i~przeglądarki \textit{Chromium}. Testowany jest cały system w~takiej postaci, w~jakiej jest wdrażany: baza danych,
emulator \textit{Azure Blob Storage}, API oraz aplikacja kliencka serwowana przez \textit{nginx}. Wszystkie te usługi
uruchamia \textit{Docker Compose} na podstawie pliku \texttt{compose.yaml}, uzupełnionego o~plik
\texttt{compose.e2e.yaml}. Plik ten przenosi dane bazy i~repozytoriów do pamięci (\textit{tmpfs}), więc każde
uruchomienie testów zaczyna się od pustego systemu. Usługi startuje i~zatrzymuje sam \textit{Playwright}. Fragment
konfiguracji przedstawiono na listingu~\ref{lst:e2e-konfiguracja}.

\begin{lstlisting}[basicstyle=\footnotesize\ttfamily, captionpos=t, caption={Fragment konfiguracji \textit{Playwright} uruchamiającej system przed testami}, label={lst:e2e-konfiguracja}]
webServer: {
  command:
    'docker compose -p projectcowork-e2e -f compose.yaml -f compose.e2e.yaml up --build --abort-on-container-failure',
  cwd: '../..',
  env: { API_PORT: '18080', APP_PORT: '13000' },
  url: `${API_URL}/api/permissions`,
  reuseExistingServer: !process.env.CI,
  timeout: 15 * 60_000,
},
\end{lstlisting}

Każdy test pracuje na własnym, nowo zarejestrowanym koncie. Testy klikają w~interfejsie tylko te elementy, które
sprawdzają. Dane potrzebne do rozpoczęcia scenariusza, takie jak konta, projekty czy sprinty, są tworzone przez API,
co skraca czas testów. Logowanie również odbywa się przez API za pomocą klienta \texttt{page.request}, który dzieli
ciasteczka z~przeglądarką. Dzięki temu ciasteczko z~tokenem odświeżania trafia do przeglądarki, a~aplikacja kliencka
odtwarza z~niego sesję tak samo jak po zalogowaniu przez formularz. Mechanizm ten przedstawiono na listingu~\ref{lst:e2e-logowanie}.

\begin{lstlisting}[basicstyle=\footnotesize\ttfamily, captionpos=t, caption={Logowanie użytkownika testowego przez API}, label={lst:e2e-logowanie}]
export async function signIn(page: Page, account: Account) {
  const response = await page.request.post(`${API_URL}/api/auth/login`, {
    data: { email: account.email, password: account.password },
  });
  expect(response.status(), await response.text()).toBe(200);
}

export const test = base.extend<Fixtures>({
  user: async ({}, use) => {
    await use(await registerAccount('user'));
  },
  userPage: async ({ page, user }, use) => {
    await signIn(page, user);
    await use(page);
  },
});
\end{lstlisting}

Elementy interfejsu są wyszukiwane przede wszystkim po roli i~nazwie dostępnej, np. przycisk o~nazwie
\textit{Zaloguj się} albo okno dialogowe o~tytule \textit{Nowe zadanie}. Takie selektory odpowiadają temu, co widzi
użytkownik, i~nie zależą od struktury kodu HTML. W~kilku miejscach, w~których element nie ma unikalnej nazwy
(karty zadań i~sekcje sprintów), dodano do komponentów atrybut \texttt{data-testid}.

\subsection{Scenariusz 1: rejestracja nowego użytkownika}
Scenariusz odtwarza pierwsze wejście użytkownika do systemu. Użytkownik zakłada konto w~formularzu rejestracji, a~system
od razu go loguje i~otwiera pulpit. Opis scenariusza zawiera tabela~\ref{tab:e2e-rejestracja}.

\begin{table}[htb]
    \centering
    \caption{Scenariusz testowy: rejestracja nowego użytkownika}
    \label{tab:e2e-rejestracja}
    \small
    \begin{tabularx}{\linewidth}{@{}>{\raggedright\arraybackslash}p{3.5cm} >{\raggedright\arraybackslash}X@{}}
        \hline
        Test & \texttt{auth.spec.ts}: \textit{rejestracja przez formularz loguje i~otwiera pulpit} \\ \hline
        Warunki początkowe & Użytkownik niezalogowany na stronie \texttt{/auth} \\ \hline
        Dane wejściowe & Unikalna nazwa użytkownika, adres e-mail i~hasło \texttt{Test1234!} \\ \hline
        Kroki & 1. Wybranie zakładki \textit{Rejestracja}. \newline 2. Wypełnienie pól \textit{Nazwa użytkownika}, \textit{E-mail}, \textit{Hasło} i~\textit{Powtórz hasło}. \newline 3. Kliknięcie przycisku \textit{Zarejestruj się} \\ \hline
        Oczekiwany wynik & Przejście na stronę główną \texttt{/} z~nagłówkiem \textit{Pulpit} \\ \hline
    \end{tabularx}
\end{table}

Implementację testu przedstawiono na listingu~\ref{lst:e2e-rejestracja}. Pole hasła jest wyszukiwane po dokładnej
nazwie, ponieważ nazwa \textit{Hasło} jest też częścią nazwy pola \textit{Powtórz hasło}.

\begin{lstlisting}[basicstyle=\footnotesize\ttfamily, captionpos=t, caption={Test rejestracji przez formularz}, label={lst:e2e-rejestracja}]
test('rejestracja przez formularz loguje i otwiera pulpit', async ({ page }) => {
  const username = unique('ui');

  await page.goto('/auth');
  await page.getByRole('tab', { name: 'Rejestracja' }).click();
  const panel = page.getByRole('tabpanel', { name: 'Rejestracja' });

  await panel.getByLabel('Nazwa użytkownika').fill(username);
  await panel.getByLabel('E-mail').fill(`${username}@example.com`);
  await panel.getByRole('textbox', { name: 'Hasło', exact: true }).fill(PASSWORD);
  await panel.getByLabel('Powtórz hasło').fill(PASSWORD);
  await panel.getByRole('button', { name: 'Zarejestruj się' }).click();

  await expect(page).toHaveURL(/\/$/);
  await expect(page.getByRole('heading', { name: 'Pulpit' })).toBeVisible();
});
\end{lstlisting}

\subsection{Scenariusz 2: wylogowanie we wszystkich kartach}
Użytkownik może mieć aplikację otwartą w~kilku kartach przeglądarki. Wylogowanie w~jednej z~nich musi zakończyć sesję
we wszystkich, żeby na współdzielonym komputerze nie została otwarta karta z~dostępem do konta. Karty powiadamiają się
o~wylogowaniu przez mechanizm \texttt{BroadcastChannel} przeglądarki. Opis scenariusza zawiera tabela~\ref{tab:e2e-karty}.

\begin{table}[htb]
    \centering
    \caption{Scenariusz testowy: wylogowanie we wszystkich kartach}
    \label{tab:e2e-karty}
    \small
    \begin{tabularx}{\linewidth}{@{}>{\raggedright\arraybackslash}p{3.5cm} >{\raggedright\arraybackslash}X@{}}
        \hline
        Test & \texttt{auth.spec.ts}: \textit{wylogowanie w~jednej karcie wylogowuje drugą} \\ \hline
        Warunki początkowe & Zalogowany użytkownik z~dwiema kartami otwartymi na stronie \texttt{/projects} \\ \hline
        Kroki & Kliknięcie przycisku \textit{Wyloguj} w~pierwszej karcie \\ \hline
        Oczekiwany wynik & Druga karta, bez żadnej akcji użytkownika, przechodzi na stronę logowania \texttt{/auth} \\ \hline
    \end{tabularx}
\end{table}

Implementację testu przedstawiono na listingu~\ref{lst:e2e-karty}. Obie karty należą do tego samego kontekstu
przeglądarki, więc dzielą ciasteczka tak jak karty w~zwykłej przeglądarce.

\begin{lstlisting}[basicstyle=\footnotesize\ttfamily, captionpos=t, caption={Test wylogowania w wielu kartach}, label={lst:e2e-karty}]
test('wylogowanie w jednej karcie wylogowuje drugą', async ({ userPage: page }) => {
  const secondTab = await page.context().newPage();
  await page.goto('/projects');
  await secondTab.goto('/projects');
  await expect(secondTab.getByRole('heading', { name: 'Projekty' })).toBeVisible();

  await page.getByRole('button', { name: 'Wyloguj' }).click();

  await expect(secondTab).toHaveURL(/\/auth$/);
});
\end{lstlisting}

\subsection{Scenariusz 3: zaproszenie użytkownika do projektu}
Scenariusz obejmuje dwóch użytkowników, z~których każdy korzysta z~własnej przeglądarki. Właściciel projektu zaprasza
drugiego użytkownika, a~ten dowiaduje się o~zaproszeniu z~ikony dzwonka w~nagłówku i~dołącza do projektu.
Na końcu właściciel widzi nowego członka na liście. Opis scenariusza zawiera tabela~\ref{tab:e2e-zaproszenie}.

\begin{table}[htb]
    \centering
    \caption{Scenariusz testowy: zaproszenie użytkownika do projektu}
    \label{tab:e2e-zaproszenie}
    \small
    \begin{tabularx}{\linewidth}{@{}>{\raggedright\arraybackslash}p{3.5cm} >{\raggedright\arraybackslash}X@{}}
        \hline
        Test & \texttt{projects.spec.ts}: \textit{zaproszenie przez UI, akceptacja dzwonkiem i~nowy członek na liście} \\ \hline
        Warunki początkowe & Dwóch zarejestrowanych użytkowników. Pierwszy z~nich jest właścicielem projektu \\ \hline
        Kroki & 1. Właściciel otwiera zakładkę \textit{Członkowie i~role}, klika \textit{Zaproś}, wpisuje adres e-mail drugiego użytkownika i~zatwierdza. \newline 2. Zaproszony klika ikonę dzwonka, a~następnie przycisk \textit{Dołącz}. \newline 3. Właściciel odświeża stronę \\ \hline
        Oczekiwany wynik & Zaproszony trafia na stronę projektu. Na liście członków u~właściciela widnieje nowy użytkownik z~rolą domyślną \textit{Członek} \\ \hline
    \end{tabularx}
\end{table}

Implementację testu przedstawiono na listingu~\ref{lst:e2e-zaproszenie}. Funkcja \texttt{openAs} tworzy osobny
kontekst przeglądarki, czyli osobny zestaw ciasteczek, i~loguje w~nim wskazanego użytkownika.

\begin{lstlisting}[basicstyle=\footnotesize\ttfamily, captionpos=t, caption={Test zaproszenia z udziałem dwóch użytkowników}, label={lst:e2e-zaproszenie}]
test('zaproszenie przez UI, akceptacja dzwonkiem i nowy członek na liście', async ({ browser, user: owner }) => {
  const invitee = await registerAccount();
  const project = await createProject(owner);

  const ownerPage = await openAs(browser, owner);
  await ownerPage.goto(`/project/${project.id}?tab=access`);
  await ownerPage.getByRole('button', { name: 'Zaproś' }).click();
  const dialog = ownerPage.getByRole('dialog', { name: 'Zaproś do projektu' });
  await dialog.getByLabel('E-mail').fill(invitee.email);
  await dialog.getByRole('button', { name: 'Zaproś' }).click();
  await expect(dialog).toBeHidden();

  const inviteePage = await openAs(browser, invitee);
  await inviteePage.goto('/');
  await inviteePage.getByRole('button', { name: 'Masz zaproszenie do projektu' }).click();

  await expect(inviteePage).toHaveURL(/\/account#invitations$/);
  await expect(inviteePage.getByText(project.name)).toBeVisible();
  await inviteePage.getByRole('button', { name: 'Dołącz' }).click();

  await expect(inviteePage).toHaveURL(new RegExp(`/project/${project.id}$`));
  await expect(inviteePage.getByRole('heading', { name: project.name, level: 1 })).toBeVisible();

  await ownerPage.reload();
  await expect(ownerPage.getByText(invitee.username, { exact: true })).toBeVisible();
  await expect(ownerPage.getByRole('combobox', { name: `Rola użytkownika ${invitee.username}` }))
    .toHaveValue('Członek');
});
\end{lstlisting}

\subsection{Scenariusz 4: zadanie w~sprincie i~zmiana jego statusu}
Scenariusz odtwarza podstawową pracę zespołu ze sprintem. Użytkownik dodaje zadanie do sprintu, otwiera stronę sprintu
i~oznacza zadanie jako zrobione. Test sprawdza, czy zmiana jest zapisana na serwerze, czyli czy zostaje po odświeżeniu
strony, oraz czy podsumowanie sprintu na stronie projektu uwzględnia ukończone zadanie. Opis scenariusza zawiera tabela~\ref{tab:e2e-sprint}.

\begin{table}[htb]
    \centering
    \caption{Scenariusz testowy: zadanie w~sprincie i~zmiana jego statusu}
    \label{tab:e2e-sprint}
    \small
    \begin{tabularx}{\linewidth}{@{}>{\raggedright\arraybackslash}p{3.5cm} >{\raggedright\arraybackslash}X@{}}
        \hline
        Test & \texttt{management.spec.ts}: \textit{nowe zadanie w~sprincie, zmiana statusu na \quotedblbase{}Zrobione\textquotedblright{}} \\ \hline
        Warunki początkowe & Zalogowany właściciel projektu z~jednym sprintem \\ \hline
        Kroki & 1. W~zakładce \textit{Sprinty i~zadania} kliknięcie \textit{Nowe zadanie}, wpisanie nazwy, wybór priorytetu \textit{Pilny} i~sprintu, zatwierdzenie. \newline 2. Przejście na stronę sprintu i~zmiana statusu zadania na \textit{Zrobione}. \newline 3. Odświeżenie strony. \newline 4. Powrót na stronę projektu \\ \hline
        Oczekiwany wynik & 1. Sekcja sprintu pokazuje \textit{1 zadanie} i~\textit{Ukończone: 0/1}. \newline 2. Zadanie ma numer \#1 i~datę ukończenia. \newline 3. Status \textit{Zrobione} pozostaje. \newline 4. Sekcja sprintu pokazuje \textit{Ukończone: 1/1} \\ \hline
    \end{tabularx}
\end{table}

Implementację testu przedstawiono na listingu~\ref{lst:e2e-sprint}. Pola wyboru biblioteki \textit{Mantine} mają rolę
\texttt{combobox}, a~lista opcji jest wyświetlana poza oknem dialogowym, dlatego funkcja \texttt{pick} najpierw klika
pole, a~następnie opcję na całej stronie.

\begin{lstlisting}[basicstyle=\footnotesize\ttfamily, captionpos=t, caption={Test dodania zadania do sprintu i zmiany jego statusu}, label={lst:e2e-sprint}]
test('nowe zadanie w sprincie, zmiana statusu na „Zrobione”', async ({ userPage: page, user }) => {
  const project = await createProject(user);
  const sprint = await createSprint(user, project.id);
  const taskName = unique('Zadanie');

  await page.goto(`/project/${project.id}?tab=sprints`);
  await page.getByRole('button', { name: 'Nowe zadanie' }).click();
  const dialog = page.getByRole('dialog', { name: 'Nowe zadanie' });
  await dialog.getByLabel('Nazwa').fill(taskName);
  await pick(page, dialog.getByRole('combobox', { name: 'Priorytet' }), 'Pilny');
  await pick(page, dialog.getByRole('combobox', { name: 'Sprint' }), sprint.name);
  await dialog.getByRole('button', { name: 'Utwórz' }).click();
  await expect(dialog).toBeHidden();

  const section = sprintSection(page, sprint.name);
  await expect(section.getByText('1 zadanie')).toBeVisible();
  await expect(section.getByText('Ukończone: 0/1')).toBeVisible();

  await section.getByRole('link', { name: 'Otwórz sprint' }).click();
  const task = taskCard(page, taskName);
  await expect(task.getByText('#1')).toBeVisible();

  await pick(page, task.getByRole('combobox', { name: 'Status' }), 'Zrobione');
  await expect(task.getByText(/^Ukończono: \d{4}-\d{2}-\d{2}$/)).toBeVisible();

  await page.reload();
  await expect(taskCard(page, taskName).getByRole('combobox', { name: 'Status' })).toHaveValue('Zrobione');

  await page.goto(`/project/${project.id}?tab=sprints`);
  await expect(sprintSection(page, sprint.name).getByText('Ukończone: 1/1')).toBeVisible();
});
\end{lstlisting}

\subsection{Scenariusz 5: przeglądanie kodu repozytorium}
Przed testem do repozytorium projektu wypychane są pliki poleceniem \texttt{git push}, tak jak zrobiłby to programista. Następnie użytkownik przegląda
w~aplikacji strukturę katalogów, otwiera plik i~sprawdza historię commitów. Opis scenariusza zawiera tabela~\ref{tab:e2e-repozytorium}.

\begin{table}[htb]
    \centering
    \caption{Scenariusz testowy: przeglądanie kodu repozytorium}
    \label{tab:e2e-repozytorium}
    \small
    \begin{tabularx}{\linewidth}{@{}>{\raggedright\arraybackslash}p{3.5cm} >{\raggedright\arraybackslash}X@{}}
        \hline
        Test & \texttt{repositories.spec.ts}: \textit{po git push widać pliki, foldery, zawartość i~commity} \\ \hline
        Warunki początkowe & Repozytorium projektu, do którego wypchnięto jeden commit z~plikami \texttt{README.md} i~\texttt{src/main.ts} \\ \hline
        Kroki & 1. Otwarcie repozytorium ze strony projektu. \newline 2. Kliknięcie folderu \texttt{src}, a~następnie pliku \texttt{main.ts}. \newline 3. Przejście do zakładki \textit{Commity} \\ \hline
        Oczekiwany wynik & 1. Folder \texttt{src} jest wyświetlony przed plikiem \texttt{README.md}. \newline 2. Adres strony zawiera ścieżkę folderu i~pliku, a~na stronie widać zawartość pliku. \newline 3. Lista zawiera commit \textit{Pierwszy commit z~E2E} autorstwa \textit{E2E Tester} \\ \hline
    \end{tabularx}
\end{table}

Implementację testu przedstawiono na listingu~\ref{lst:e2e-repozytorium}. Funkcja \texttt{pushFiles} tworzy
w~katalogu tymczasowym lokalne repozytorium z~podanymi plikami i~wypycha je na adres zwrócony przez API.

\begin{lstlisting}[basicstyle=\footnotesize\ttfamily, captionpos=t, caption={Test przeglądania kodu i historii repozytorium}, label={lst:e2e-repozytorium}]
test('po git push widać pliki, foldery, zawartość i commity', async ({ userPage: page, user }) => {
  const project = await createProject(user);
  const repository = await createRepository(user, project.id);
  pushFiles(
    repository.cloneUrl,
    { 'README.md': '# Projekt E2E\n', 'src/main.ts': 'export const answer = 42;\n' },
    'Pierwszy commit z E2E',
  );

  await page.goto(`/project/${project.id}`);
  await page.getByRole('link', { name: `Otwórz repozytorium ${repository.name}` }).click();

  const files = page.getByRole('table');
  await expect(files.getByRole('row').nth(0)).toContainText('src');
  await expect(files.getByRole('row').nth(1)).toContainText('README.md');

  await files.getByText('src', { exact: true }).click();
  await expect(page).toHaveURL(/path=src/);
  await files.getByText('main.ts').click();

  await expect(page).toHaveURL(/file=src%2Fmain\.ts|file=src\/main\.ts/);
  await expect(page.getByText('export const answer = 42;')).toBeVisible();

  await page.getByRole('tab', { name: 'Commity' }).click();
  await expect(page.getByText('Pierwszy commit z E2E')).toBeVisible();
  await expect(page.getByText(/^E2E Tester · /)).toBeVisible();
});
\end{lstlisting}